\section{Variance}

\subsection{Non-Equal Likelihood}

The {\bf Variance} of $X$ is defined as 
\begin{align}
\var(X) \; \defe \; & \; 
	\E\left[ (X - \ex)^2 \right], 
	\label{eq:var_definition} \\
 & \; \E\left[ (X - \ex) \; (X - \ex) \right], \nonumber \\ 
 & \; \E\left[ X^2 - 2 \; \ex \; X + \ex \; \ex \right], \nonumber \\
 & \; \E[X^2] - 2 \; \ex \; \ex \;  + \ex \; \ex,  \nonumber \\
 & \; \E[X^2] - (\ex)^2. \label{eq:variance}
\end{align}

\subsection{Equal Likelihood}
For the special case where each probability is equally likely, 
(\ref{eq:variance}) takes the form
\be
\var(X)
\overset{(\ref{eq:variance})}{=} 
\sum_{i=1}^n \; p_i x^2_i - \mu^2 = \frac{1}{n} \sum_{i=1}^n \; x^2_i - \mu^2 
\defe
\sigma_x^2,
\label{eq:simple_variance}
\ee
and is referred to as the {\bf simple variance}.

\begin{Hremark}
The form of (\ref{eq:simple_variance}) is often stated in slightly
different form, 
\begin{align}
\sigma_x^2 = 
\frac{1}{n} \sum_{i=1}^n \; x^2_i - \mu^2 & =
\frac{1}{n} \sum_{i=1}^n \; x^2_i - 2 \mu^2 + \mu^2, \\
& \overset{(\ref{eq:simple_average})}{=} \frac{1}{n} \sum_{i=1}^n \; x^2_i -
    \frac{2\mu}{n} \sum_{i=1}^n \; x_i +
	\frac{n \mu^2}{n}, \\
& = \frac{1}{n} \sum_{i=1}^n \left( x^2_i - 2 \mu x_i + \mu^2 \right), \\
& = \frac{1}{n} \sum_{i=1}^n \left( x_i - \mu \right)^2.
\end{align}
\end{Hremark}

\section{Covariance}
\subsection{Non-Equal Likelihood}

For two jointly distributed real-valued random variables $X$ and $Y$, with
{\bf expected values}\footnote{Also known as {\bf mean values}.} 
$\ex$ and $\ey$, respectively, the {\bf covariance} is defined as
\begin{align}
\cov(X,Y) \; \defe \; & \; \E\left[ 
(X - \ex)) \; (X - \ey))
\right], \label{eq:cov_definition} \\
 & \; \exy - \ex \; \ey - \ex \; \ey + \ex \; \ey, \nonumber \\
 & \; \exy - \ex \; \ey.
\end{align}

Let the random variable pair $(X,Y)$ take on values $(x_i, y_i)$ for
$i=1,\ldots,n$ with probability $p_i$.  The covariance is defined as
\be
\cov(X,Y) \defe \sum_{i=1}^n \; p_i \; \left(
x_i - \ex
\right)
\;
\left(
y_i - \ey
\right).
\label{eq:cov_discrete}
\ee

\begin{Hremark}
The covariance of a random variable with itself is its 
variance, 
\be
\cov(X,X) = \var(X), 
\ee
which is seen by comparing 
(\ref{eq:cov_definition})
to 
(\ref{eq:var_definition}).
\end{Hremark}

\subsection{Equal Likelihood}
If the probability of each outcome is equal, the $p_i = 1/n$, and
(\ref{eq:cov_discrete}) becomes
\be
\cov(X,Y) = \frac{1}{n} \; \sum_{i=1}^n \; \left(
x_i - \mu_x
\right)
\;
\left(
y_i - \mu_y
\right).
\label{eq:cov_discrete_equal}
\ee

\section{Correlation}

\subsection{Pearson Correlation Coefficient}
The {\bf Pearson correlation coefficient} (PCC)\footnote{Also known as the
Pearson's $r$ coefficient, Pearson product-moment correlation
coefficient (PPMCC), and the bivariate correlation.} is a 
measure of the linear correlation between two variables 
$X$ and $Y$.  

When a pair of random variables $(X,Y)$ are applied to a population, the PCC
is called the {\bf population correlation coefficient} and denoted as 
$\rho_{X,Y}$ and is defined as 
\be
\rho_{X,Y} \defe \frac{\cov(X,Y)}{\sigma_X \; \sigma_Y}, 
\ee
where 
\begin{itemize}
\item $\cov$ is the covariance of $X$ and $Y$, 
\item $\sigma_X$ is the standard deviation of $X$, and 
\item $\sigma_Y$ is the standard deviation of $Y$.
\end{itemize}
The PCC is bounded as
\be
-1 \leq \rho_{X,Y} \leq 1.
\ee
\begin{itemize}
\item A value of -1 indicates complete negative linear correlation.
\item A value of  0 indicates the absence of linear correlation.
\item A value of +1 indicates complete positive linear correlation.
\end{itemize}

\begin{Hremark}
	When $\exy \approx \ex \; \ey$, 
	then the $\cov(X,Y) = \exy - \ex \; \ey$ 
	is prone to so-called ``catastrophic 
	cancellation'' with floating point
	calculations.  A numerically stable algorithm 
	should be used.
\end{Hremark}
